\documentclass[11pt,a4paper]{article}
\usepackage{geometry}
\geometry{margin=1in}

\title{\textbf{Data Dictionary}\\[0.5cm] \textbf{AeroGrid}\\[0.3cm] Low-Altitude Airspace Intelligence Warehouse}
\author{AeroGrid Development Team}
\date{\today}

\begin{document}

\maketitle

\section{Introduction}
This document serves as the formal Data Dictionary for the AeroGrid Data Warehouse, detailing all dimension and fact tables, columns, data types, key classifications, and business descriptions.

\section{Summary of Tables}
\begin{itemize}
    \item \textbf{Dimensions:} \texttt{DimTime}, \texttt{DimVehicle}, \texttt{DimAirspaceZone}, \texttt{DimWeather}, \texttt{DimSpatialGrid}.
    \item \textbf{Facts:} \texttt{FactVehicleTracking}, \texttt{FactFlightJourney}, \texttt{FactCloseEncounter}, \texttt{FactRestrictionActivation}, \texttt{FactAirspaceOccupancy}.
\end{itemize}

\section{Key Standards}
\begin{itemize}
    \item Surrogate keys (\texttt{\_key}) are utilized across all dimension tables.
    \item Natural keys (\texttt{\_id}) preserve source system identifiers for lineage tracking.
    \item Altitude measures are standardized to feet MSL (\texttt{altitude\_ft}).
\end{itemize}

\section{Detailed Schema Reference}
For complete column definitions, data types, and analytical use cases, refer to the master data dictionary in \texttt{data/data\_dictionary.md}.

\end{document}
